\section*{Bab \ref{ch:g11:absorbance} --- Warna dan Absorbans}

\begin{solution}{exo:g11:absorbance:1}
Karena menyerap jingga, larutan itu tampak biru, warna yang berhadapan
dengan jingga pada roda. Karena menyerap ungu, larutan itu tampak
kuning.
\end{solution}

\begin{solution}{exo:g11:absorbance:2}
Larutan hijau terutama menyerap cahaya merah, warna yang berhadapan
dengan hijau pada roda: dari sekitar \qty{620}{nm} sampai \qty{750}{nm}.
\end{solution}

\begin{solution}{exo:g11:absorbance:3}
Zat warna biru: $\lambda_{\max} = \qty{629}{nm}$, $A = 0.82$. Zat warna
kuning: $\lambda_{\max} = \qty{427}{nm}$, $A = 0.795$, sekitar 0.80.
\end{solution}

\begin{solution}{exo:g11:absorbance:4}
$A$ sebanding dengan konsentrasi: tiga kali lebih banyak memberikan
$A = 0.90$; setengahnya memberikan $A = 0.15$.
\end{solution}

\begin{solution}{exo:g11:absorbance:5}
Blanko adalah kuvet yang berisi \omterm{def:g6:solutions-and-solubility:solute}{pelarut} murni. Menolkan $A = 0$ dengannya
menghapus apa yang diserap kuvet, \omterm{def:g6:solutions-and-solubility:solute}{pelarut}, dan alat itu, sehingga
pembacaan-pembacaan sesudahnya hanya disebabkan oleh \omterm{def:g6:solutions-and-solubility:solute}{zat terlarut}.
\end{solution}

\begin{solution}{exo:g11:absorbance:6}
$C = 0.41 / 0.164 = \qty{2.5}{mg/L}$.
\end{solution}

\begin{solution}{exo:g11:absorbance:7}
Pada \qty{629}{nm} \omterm{def:g11:absorbance:absorbance}{absorbansnya} paling besar, sehingga pengukurannya
paling peka (konsentrasi yang kecil masih memberikan \omterm{def:g11:absorbance:absorbance}{absorbans} yang
dapat dibaca); dan di puncak pita kurvanya datar, sehingga sedikit
kesalahan panjang gelombang hampir tidak mengubah pembacaannya. Pada
\qty{550}{nm} \omterm{def:g11:absorbance:absorbance}{absorbansnya} kecil dan berubah cepat menurut panjang
gelombang.
\end{solution}

\begin{solution}{exo:g11:absorbance:8}
Setelah diencerkan: $C = 0.48 / 0.164 = \qty{2.9}{mg/L}$; minumannya:
$5 \times 2.9 = \qty{15}{mg/L}$ (lebih tepatnya $14.6$).
\end{solution}

\begin{solution}{exo:g11:absorbance:9}
$a = A / (\ell\, C_m) = 0.795 / (1.00 \times 0.0150) = \qty{53.0}{L/(g.cm)}$;
$\varepsilon = a \times M = 53.0 \times 534.4 = \qty{2.83e4}{L/(mol.cm)}$.
\end{solution}

\begin{solution}{exo:g11:absorbance:10}
Di atas sekitar \qty{520}{nm}. Pada \qty{629}{nm} zat warna kuning sama
sekali tidak menyerap, sehingga seluruh \omterm{def:g11:absorbance:absorbance}{absorbans} di situ berasal dari
zat warna biru.
\end{solution}

\begin{solution}{exo:g11:absorbance:11}
\omterm{def:g11:absorbance:absorbance}{Absorbansnya} menjadi dua kali lipat: $A = \varepsilon \ell c$ sebanding
dengan $\ell$; cahaya melewati larutan dua kali lebih tebal, sehingga
melewati \omterm{def:g7:atoms-and-molecules:molecule}{molekul} penyerap dua kali lebih banyak.
\end{solution}

\begin{solution}{exo:g11:absorbance:12}
$A/C$: $0.164$, $0.162$, $0.115$, $0.0675$. Perbandingannya hanya tetap
untuk dua yang pertama: di atas sekitar \qty{10}{mg/L} (\omterm{def:g11:absorbance:absorbance}{absorbans} di
atas sekitar 1.6) hukum itu tidak berlaku lagi. Hanya \omterm{def:g10:concentration-and-dilution:standard}{larutan standar} 5
dan \qty{10}{mg/L} (dan yang lebih rendah) yang boleh dipakai; larutan
tak diketahui yang menunjukkan 2.5 harus diencerkan, misalnya lima kali,
lalu diukur lagi.
\end{solution}

\begin{solution}{exo:g11:absorbance:13}
Zat warna biru, dari \qty{629}{nm} saja:
$C = 0.574 / 0.164 = \qty{3.5}{mg/L}$. Zat warna kuning, dari
\qty{427}{nm}: $C = 0.53 / 0.0530 = \qty{10}{mg/L}$.
\end{solution}

\begin{solution}{exo:g11:absorbance:14}
Ia memperoleh $0.368 / 0.164 = \qty{2.24}{mg/L}$. \omterm{def:g11:absorbance:absorbance}{Absorbans} yang
sebenarnya $0.368 - 0.040 = 0.328$, sehingga konsentrasi yang sebenarnya
$0.328 / 0.164 = \qty{2.00}{mg/L}$. Hasilnya \qty{12}{\percent} terlalu
tinggi.
\end{solution}

\begin{solution}{exo:g11:absorbance:15}
$10 \times 60 = \qty{600}{mg}$ per hari, dalam $600 / 20 = \qty{30}{L}$
minuman. Tidak ada orang yang minum tiga puluh liter sehari: minuman itu
sendiri bukan risiko yang realistis, meskipun zat warna itu dapat pula
berasal dari makanan lain.
\end{solution}

\begin{solution}{pb:g11:absorbance:1}
\textbf{1.} Jingga kemerahan, warna yang berhadapan dengan biru pada
roda.

\textbf{2.} $\lambda_{\max} = \qty{629}{nm}$.

\textbf{3.} Pada \qty{450}{nm} cahayanya biru: zat warna itu
meloloskannya, dan itulah tepatnya sebabnya minuman itu tampak biru.

\textbf{4.} Pada \qty{629}{nm}.

\textbf{5.} $\qty{50.0}{mg} / \qty{0.5000}{L} = \qty{100}{mg/L}$.

\textbf{6.} \omterm{def:g10:concentration-and-dilution:dilution}{Faktor pengenceran} $100 / 3.0 = 33.3$, jadi
$V_0 = 100.0 / 33.3 = \qty{3.00}{mL}$: pipet ukur (atau buret) dan labu
ukur \qty{100.0}{mL}.

\textbf{7.} $A/C$ = 0.168, 0.163, 0.165, 0.163, 0.165: semuanya dekat
dengan 0.164, jadi $A \approx 0.164\, C$, sebuah garis yang melalui
titik asal.

\textbf{8.} Larutan tanpa zat warna ($C = 0$) adalah blanko, yang
dinolkan pada $A = 0$.

\textbf{9.} $a = \qty{0.164}{L/mg}$ per sentimeter, yaitu
\qty{164}{L/(g.cm)}; $\varepsilon = 164 \times 792.9 =
\qty{1.30e5}{L/(mol.cm)}$.

\textbf{10.} $F = 50.0 / 25.0 = 2$.

\textbf{11.} $C = 0.590 / 0.164 = \qty{3.60}{mg/L}$.

\textbf{12.} $2 \times 3.60 = \qty{7.20}{mg/L}$.

\textbf{13.} Sekitar $2 \times 0.590 = 1.18$, di atas \omterm{def:g10:concentration-and-dilution:standard}{larutan standar}
tertinggi (0.823): pembacaannya akan berada di luar rentang yang
dikalibrasi, tempat hukum itu mungkin tidak berlaku. \omterm{def:g10:concentration-and-dilution:dilution}{Pengenceran}
membawanya ke antara larutan-larutan standar.

\textbf{14.} $\qty{7.20}{mg/L} \times \qty{0.500}{L} = \qty{3.60}{mg}$.

\textbf{15.} $n = \num{3.60e-3} / 792.9 = \qty{4.54e-6}{mol}$.

\textbf{16.} $6 \times 30 = \qty{180}{mg}$ per hari.

\textbf{17.} $180 / 3.60 = 50$ botol.

\textbf{18.} $50 \times 0.500 = \qty{25}{L}$ dalam sehari: mustahil
diminum. Zat warna minuman ini saja tidak dapat membawa seorang anak
mendekati batas itu, walaupun zat warna dari berbagai makanan saling
menjumlah.

\textbf{19.} $6 \times 70 = \qty{420}{mg}$, yaitu $420 / 3.60 \approx
117$ botol.

\textbf{20.} Sekitar 50 botol \qty{500}{mL}, \qty{25}{L} minuman, dalam
satu hari saja.
\end{solution}
